% Resolve the shared preamble when compiling from fundamentals/.
\makeatletter
\def\input@path{{./}{../}{../../}{../../../}{../../../../}{preamble/}{../preamble/}{../../preamble/}{../../../preamble/}}
\makeatother
\PassOptionsToClass{letterpaper}{article}
\input{preamble_exams.tex}

% Worksheet-only overrides; all packages come from the shared preamble.
\pagestyle{empty}
\linespread{1}
\setlength{\parskip}{0pt}
\setlist[enumerate]{leftmargin=*,itemsep=2pt,parsep=0pt,topsep=2pt}
\makeatletter
\renewcommand\section{\@startsection{section}{1}{0pt}{6pt}{2pt}{\normalfont\normalsize\bfseries}}
\makeatother
\newcommand{\challenge}[1]{\par\textit{Challenge: #1}}

\begin{document}
\newgeometry{left=0.4in,right=0.4in,top=0.35in,bottom=0.35in}
% Standard 10pt text keeps all twelve problems readable on one print page.
\small
\begingroup
\centering
{\large\bfseries Probability Challenge\par}
{\small Olympiad-Style Problems for Grades 7--8\par}
\endgroup
\vspace{4pt}

\textbf{Instructions.} Solve each problem using logical reasoning, diagrams, tables, or organized counting. Show your work and justify your probability calculations. Unless stated otherwise, assume all random choices are independent and equally likely among the specified possibilities.

\section*{Level 1 --- Clever Counting}
\begin{enumerate}[label=\textbf{\arabic*.}]
\item \textbf{The Two-Dice Trap.} Two fair six-sided dice are rolled. You win if their sum is greater than their product. What is the probability of winning?
\challenge{Find a solution without listing all 36 ordered outcomes individually.}

\item \textbf{The Secret Three-Digit Code.} A computer randomly generates a code consisting of three digits. Each digit is independently selected from 0 to 9, and repetitions are allowed. What is the probability that the digits are in strictly increasing order from left to right? For example, \texttt{148} qualifies, but \texttt{133} and \texttt{941} do not.
\challenge{Count the successful codes without listing them.}

\item \textbf{The Frog's Lucky Jump.} A frog starts at position 0 on a number line and makes four jumps. Each jump independently moves it one step left or one step right, with equal probability. What is the probability that the frog returns to position 0 after its fourth jump?
\challenge{Are all possible final positions equally likely?}

\item \textbf{The Birthday Calendar.} Two students independently choose a random day of the week, from Monday through Sunday. Each day is equally likely. What is the probability that their chosen days are either the same or consecutive? Sunday and Monday are considered consecutive.
\challenge{Explain why the circular arrangement of the week matters.}
\end{enumerate}

\section*{Level 2 --- Strategic Reasoning}
\begin{enumerate}[label=\textbf{\arabic*.},start=5]
\item \textbf{The Last Candy.} A bag contains 3 red candies and 3 blue candies. Candies are removed one at a time, uniformly at random, without replacement. What is the probability that the last candy is red, given that the first two candies removed have different colors?
\challenge{Can you avoid constructing a complete probability tree?}

\item \textbf{The Triangle Game.} Three different numbers are selected uniformly at random from 1 through 9. What is the probability that they can be the side lengths of a triangle? Recall that the sum of the two shorter sides must be greater than the longest side.
\challenge{Count the successful triples systematically, without omissions or repetitions.}

\item \textbf{The Race to Six.} Two players independently roll a fair six-sided die. The player with the higher number wins; equal numbers result in a tie. Given that at least one player rolled a 6, what is the probability that the game ends in a tie?
\challenge{Explain why the answer is not $1/6$.}

\item \textbf{The Robot's Return.} A robot begins at the bottom-left corner of a $2\times2$ square grid. It makes exactly four moves. Each move independently goes one unit up, down, left, or right, each with probability $1/4$. The robot may leave the original grid. What is the probability that it returns to its starting position?
\challenge{Organize the possible routes by their combinations of directions.}
\end{enumerate}

\section*{Level 3 --- Olympiad Challenges}
\begin{enumerate}[label=\textbf{\arabic*.},start=9]
\item \textbf{The First Double.} A fair coin is tossed repeatedly until two consecutive tosses show the same result. What is the probability that the game ends with two heads rather than two tails?
\challenge{Solve the problem without listing infinitely many cases.}

\item \textbf{The Three Students' Hats.} Three students stand in a line, facing forward. Each receives a hat independently chosen to be red or blue, with equal probability. The student at the back can see the other two hats. The middle student can see only the front student's hat. The front student sees no hats. Before receiving the hats, the students may agree on a strategy. Afterward, each must simultaneously guess the color of their own hat, without communicating. What is the greatest possible probability that \emph{at least two students guess correctly}?
\challenge{Can a strategy based on what students see improve on independent random guessing?}

\item \textbf{The Broken Stick.} A stick of length 1 is broken at two points chosen independently and uniformly at random along its length. What is the probability that the three resulting pieces can form a triangle?
\challenge{Represent the two breaking points in a unit square and identify the successful region.}

\item \textbf{The Game of Increasing Numbers.} A player rolls a fair six-sided die three times and wins if the results are in strictly increasing order. For example, $(1,3,5)$ wins, but $(2,2,6)$ and $(4,3,6)$ lose. What is the probability of winning?

Now suppose the player rolls four times and wins if \emph{some three consecutive results} are strictly increasing. Does the probability of winning double? Determine the exact new probability.
\challenge{Avoid double-counting sequences that satisfy both winning conditions.}
\end{enumerate}

\end{document}
